\chapter{関連研究}
本章では，本研究に関連する先行研究として，推し活・ファン行動に関する研究，推し活とウェルビーイング・負担に関する研究，性格特性と個人化に関する研究，パーソナルインフォマティクスと行動変容に関する研究，アバターと将来の自己に関する研究，予測の不確かさの提示と推薦の説明に関する研究，生理指標によるストレス計測に関する研究を概観し，本研究の位置付けを述べる．

\section{推し活・ファン行動に関する研究}
ファン行動は，社会学やメディア研究の分野で古くから扱われてきた．Sandvoss\cite{sandvoss2005}は，ファンが対象に向ける熱中を，消費を通じて自己を映し出す「鏡」として捉え，ファンであることが個人のアイデンティティと深く結びついていることを論じた．

日本のファン心理については，吉田\cite{yoshida2010}が自己愛傾向とファン行動の関連を調べ，ファン行動の質問項目に基づくクラスタリングにより，自己愛傾向がファン行動のあり方に関わることを示した．小城\cite{koshiro2018}は，ファン心理尺度を再検討し，ファンが対象に向ける心理を多面的に測定する尺度を提案している．これらの研究は，ファン行動に個人差があり，その背景に心理的な傾向があることを示しているが，個人を支援することは目的としていない．

近年は「推し」や「推し活」という言葉そのものを対象とした研究も増えている．市村\cite{ichimura2025}は，女子大学生を対象に「ファン活動」「オタク活動」「推し活」のイメージと推し活を行う理由を調べ，推し活が肯定的なイメージをもつ活動として捉えられていることを示した．本間ら\cite{homma2025}は，「推し」がどのように形成され，何がそれを決めるのかをモデル化することを試みている．東ら\cite{higashi2024}は，「推し」を認知科学の観点から捉える応援研究の進展と可能性を論じている．これらの研究は推し活という現象の理解を深めているが，推し活の負担や，その個人に合わせた支援までは扱っていない．

\section{推し活とウェルビーイング・負担に関する研究}
推し活が幸福感にもたらす効果については，いくつかの定量的な検討がある．井上・上田\cite{inoue2023}は，アイドルに対するファンの心理的所有感が幸福感と関連すること，また他のファンへの意識（競争的な意識）が応援の継続を妨げうることを示した．上田・井上\cite{ueda2023}は，推し活意識が推しへの心理的所有感を媒介して幸福感を高めることを報告している．宮田\cite{miyata2024}は，幸福度と推し活の関係について考察している．

幸福やウェルビーイングの捉え方としては，人生への満足や主観的な幸福感を中心とする立場と，人生の意味や成長，他者とのよい関係などを含む心理的ウェルビーイングを重視する立場がある．西田\cite{nishida2000}は，成人女性を対象に心理的well-beingを多面的に測定する尺度を作成している．大野\cite{ono1984}は，青年の「充実感」を測定する尺度を作成し，生活への張りや主体性といった側面を捉えている．主観的幸福感については日本版主観的幸福感尺度（SHS）\cite{shimai2004}が，精神的健康の簡便な指標としてはWHO-5\cite{topp2015}とその日本語版\cite{inagaki2013}が広く用いられている．自己決定理論\cite{ryan2000}は，自律性・有能感・関係性の欲求が満たされることがウェルビーイングを支えるとしており，推し活において本人が自分で選ぶことや仲間とのつながりが重要であることとも整合する．

一方，熱心な関わりが疲弊を生むことは，対人援助職などのバーンアウト研究で知られている．Maslach と Jackson\cite{maslach1981}は，情緒的消耗感などからなるバーンアウトの測定尺度を提案し，久保\cite{kubo2007}は日本のヒューマンサービス職におけるバーンアウトを整理している．推し活における負担は職業上のバーンアウトとは異なるが，対象への強い関与が疲れを生むという点で共通する．しかし，推し活の便益と負担を同じ人の中で同時に測定し，それらがどの行動から生じるのかを分析した研究はほとんどない．

\section{性格特性と個人化に関する研究}
性格特性の代表的な枠組みとして，外向性・協調性・誠実性・神経症傾向・開放性の5因子からなるビッグファイブがある．Goslingら\cite{gosling2003}は，各因子を2項目，計10項目で測定する短縮版尺度TIPIを開発し，小塩ら\cite{oshio2012}はその日本語版TIPI-Jを作成して信頼性と妥当性を検討した．TIPI-Jは回答の負担が小さいため，アプリの中で利用者の性格を把握する用途に適している．

推薦システムや健康支援の分野では，利用者の性格に合わせて推薦や介入を個人化する試みがある．しかし，性格をどのように個人化に用いるべきか，すなわち性格を固定的な「タイプ」として扱い，タイプごとに決まった推薦を行うのか，それとも性格が行動や関わり方とどう結びついているかを踏まえて行動と組み合わせて用いるのかは，対象領域ごとに実証的に確かめる必要がある．推し活においてこの点を検討した研究はない．

\section{パーソナルインフォマティクスと行動変容に関する研究}
自分に関するデータを収集し，振り返ることで自己理解や行動変容につなげる取り組みは，パーソナルインフォマティクスと呼ばれる．Liら\cite{li2010}は，パーソナルインフォマティクスシステムの利用を「準備・収集・統合・振り返り・行動」の段階からなるモデルで整理し，各段階で利用者がつまずきうる障壁を示した．Choeら\cite{choe2014}は，自分のデータを記録・探索する人々の実践を調べ，データを集めても振り返りや行動につなげることが難しい場合が多いことを報告している．すなわち，記録を見せるだけでは行動は変わりにくい．

行動変容の理論として，Fogg\cite{fogg2009}は，行動は動機・能力・きっかけ（プロンプト）の3要素がそろったときに起こるとする行動モデルを提案し，小さく実行しやすい行動を適切なきっかけとともに示すことの重要性を説いた．Lallyら\cite{lally2010}は，新しい行動が習慣として定着するまでの期間を実生活で調べ，定着までには個人差が大きく，数週間から数か月を要することを示した．健康行動の理論については松本\cite{matsumoto2002}が，行動変容のステージや自己効力感などの概念を整理している．これらの知見は，行動変容の支援においては，本人が実行しやすい小さな行動を自分で選べるようにし，一定期間の振り返りを伴わせることが重要であることを示している．

\section{アバターと将来の自己に関する研究}
YeeとBailenson\cite{yee2007}は，仮想空間で用いるアバターの見た目が利用者自身の行動に影響を与える現象を「プロテウス効果」と名付けた．FoxとBailenson\cite{fox2009}は，本人に似たアバターが運動して体型が変化する様子を見せると，その後の運動行動が増えることを示した．Hershfieldら\cite{hershfield2011}は，年齢を重ねた将来の自分の顔を見せると，将来のための貯蓄が増えることを示した．これらの研究は，本人に似た姿で将来の状態を見せることが，現在の行動を変えるきっかけになりうることを示している．

ただし，これらの研究が対象とした運動や貯蓄は「多いほど良い」という望ましい方向が明確な領域である．推し活では，何が「良い」かは人によって異なり，活動を増やすことが良い人もいれば，一部を整えることが良い人もいる．したがって，推し活においてアバターで将来を示す場合には，一方向の理想像を押しつけるのではなく，本人が選択肢を比較して自ら選べる形にする必要がある．

\section{予測の不確かさの提示と推薦の説明に関する研究}
個人に対する予測には必ず誤差が伴う．Kayら\cite{kay2016}は，バスの到着予測の不確かさを日常のモバイルアプリで可視化する方法を検討し，一般の利用者も適切に可視化された不確かさを理解し，判断に役立てられることを示した．推し活の見込みのように誤差が大きい予測を示す場合，点の予測だけを示すと過信を招くおそれがあるため，幅とともに示すことが重要である．

推薦システムにおいては，なぜその推薦がなされたのかを説明することが，利用者の信頼や納得，意思決定の質を高めるとされる．TintarevとMasthoff\cite{tintarev2007}は，推薦の説明が果たす役割として透明性，納得性，信頼性，有効性，説得性などを整理している．本研究では，提案の理由と得点の算出式を画面に示すことで，利用者が提案を自分で評価して選べるようにする．

\section{生理指標によるストレス計測に関する研究}
ストレスの客観的な指標として，心拍変動（HRV）が用いられている．ThayerとLane\cite{thayer2007}は，迷走神経の機能を反映する心拍変動が心血管疾患や死亡のリスクと関連することを示し，自律神経の状態を捉える指標としての有用性を論じている．筆者の前研究\cite{imai2025thesis,imai2025hr}では，スマートウォッチで推し活時の心拍を測定したが，体験者が3名と少なく，推し活以外の要因（安静や移動）の影響を切り分けることも難しかった．そこで本研究では，生理指標ではなく，多数の人から得られる質問紙と週ごとの振り返りによって幸福度と推し疲れを捉え，システムによる支援に焦点を当てる．

\section{本研究の位置付け}
表\ref{tab:position}に，主な先行研究と本研究の比較を示す．先行研究は，推し活の動機や心理的な便益を記述する研究，便益の一部を定量的に検討する研究，アバターや将来の自己によって行動変容を促す研究に分かれており，推し活の便益と負担を同時に扱い，それに基づいて個人に合わせた提案を行い，その効果を検証した研究はない．筆者らは，OshiFitの構想\cite{imai2026ievc}と，同じ調査データの合成得点による初期分析\cite{imai2026g43}を報告してきた．本研究は，これらを発展させ，項目単位の分析によって設計指針を導き，システムを構築して検証するものである．

\begin{table}[t]
\centering
\jetcaption{先行研究と本研究の比較}{Comparison of prior work and this study}
\label{tab:position}
\small
\begin{tabular}{@{}lccccc@{}}
\toprule
研究 & 推し活 & 便益 & 負担 & 個人化 & 支援・検証 \\
\midrule
吉田\cite{yoshida2010}，小城\cite{koshiro2018} & ○ & △ & -- & -- & -- \\
市村\cite{ichimura2025}，本間ら\cite{homma2025} & ○ & ○ & -- & -- & -- \\
井上・上田\cite{inoue2023,ueda2023} & ○ & ○ & △ & -- & -- \\
Fox・Bailenson\cite{fox2009}，Hershfieldら\cite{hershfield2011} & -- & -- & -- & ○ & ○ \\
Choeら\cite{choe2014}，Liら\cite{li2010} & -- & -- & -- & △ & △ \\
今井ら（前研究）\cite{imai2025hr} & ○ & ○ & △ & -- & -- \\
\textbf{本研究} & ○ & ○ & ○ & ○ & ○ \\
\bottomrule
\end{tabular}

\vspace{2mm}
{\footnotesize ○：主に扱う，△：一部扱う，--：扱わない}
\end{table}
